% 第14章 空间计量
\chapter{空间计量}

\begin{chapteroverview}
\textbf{这个分类是干什么的？}

空间计量经济学研究\textbf{有地理位置信息的数据}。普通回归假设样本之间互相独立，但现实中"近朱者赤"——相邻地区的房价、污染、经济增长往往相互影响。空间计量就是把这种\textbf{空间依赖关系}量化并纳入模型，避免普通回归得出错误结论。

\textbf{什么数据需要空间分析？}

只要你的数据有\textbf{地理位置信息}（经纬度、区县、城市、邮编），就应该考虑空间分析！典型如：各区县的房价/污染/GDP/创新产出——这些指标天然存在"相邻地区互相影响"。如果忽略空间效应，回归系数和显著性都会失真。

\textbf{美赛什么题型常用？}

美赛C题（大数据分析）涉及区域经济、环境监测、城市规划时常用。Moran指数判断"污染是不是连片聚集"，空间回归量化"邻居效应有多大"。注意：\textbf{SPSS原生对空间计量支持有限}，主要通过Dabbit AI SPSS工具箱或Python的pysal库实现。

\textbf{学习路径建议}

先学\textbf{空间权重矩阵构造}（定义谁和谁是邻居）和\textbf{Moran指数}（判断有没有空间聚集），再学\textbf{空间OLS回归}（基准+空间诊断），然后是\textbf{空间滞后模型SAR}（邻居因变量影响本地）。进阶了解SEM、SDM等更复杂的空间模型。
\end{chapteroverview}

\section{基础级算法}

%% ========== 空间权重矩阵构造 ==========
\basicalgo{空间权重矩阵构造}{spatial_weights}

\begin{algointro}
空间权重矩阵W是空间分析的\textbf{地基}——它定义了"谁和谁是邻居"。没有W，后面所有空间模型都无从谈起。它是一个 $n\times n$ 矩阵，$W_{ij}$ 表示区域 $i$ 和区域 $j$ 之间的空间关联强度。
\end{algointro}

\begin{algoscene}
\begin{itemize}
    \item 任何空间分析的\textbf{第一步}——先定义邻居关系
    \item 门店网络、城市网络、区域经济分析的前置步骤
    \item 需要判断"距离越近影响越大"还是"共享边界才算邻居"
\end{itemize}
\end{algoscene}

\begin{algoprinciple}
想象你在画一张城市地图，要定义"相邻"：

\begin{itemize}
    \item \textbf{K近邻法（KNN）}：每个点找最近的K个点当邻居（如K=4）——不管实际距离多远，每个点都有4个邻居
    \item \textbf{反距离法（IDW）}：距离越近权重越大，权重 $=1/d^\alpha$，远处的点也有微弱影响
    \item \textbf{邻接法}：共享边界才算邻居（如两个区县接壤）
\end{itemize}

构造后通常做\textbf{行标准化}：每行除以行和，使权重变成"邻居的加权平均"。这样 $Wy$ 就是"邻居的平均取值"，好解释。

就像定义社交网络里"好友"的标准——是通讯录里前4个？还是距离1公里内？还是互相关注？不同定义影响后续分析。
\end{algoprinciple}

\begin{algosposs}
SPSS原生不支持空间权重矩阵，在\textbf{Dabbit AI SPSS工具箱}中操作：

\begin{enumerate}
    \item 上传含坐标的数据（x坐标列、y坐标列）
    \item 权重类型选"反距离"或"k近邻"
    \item 距离度量选"欧氏距离"（平面坐标）或"haversine"（经纬度）
    \item 近邻数K设为4，距离衰减幂设为1
    \item 标准化方式选"行标准化"
    \item 点击运行，输出连接密度、邻居数、孤立单元数等结构统计
\end{enumerate}
\end{algosposs}

\begin{algoresult}
\begin{itemize}
    \item \textbf{连接密度}：非零连接占比，1.0表示所有点都有邻居
    \item \textbf{孤立单元数}：=0最好——有孤立单元说明该点没有邻居，空间滞后无法计算
    \item \textbf{邻居数分布}：每个点平均几个邻居，是否均匀
    \item \textbf{W·y相关}：y与邻居均值的相关系数，正相关说明高值点倾向聚集
\end{itemize}
\end{algoresult}

\begin{algonotes}
\textbf{什么时候用？}任何空间分析的前置步骤。

\textbf{常见坑}
\begin{itemize}
    \item \textbf{坑1}：不做行标准化——权重直接是距离倒数，不同区域权重和差异巨大
    \item \textbf{坑2}：K设太大或太小——K=4是常用起点，可敏感性分析
\end{itemize}

\textbf{注意}：空间权重矩阵的选择会影响结论，论文中应说明选了哪种W并做敏感性检验。
\end{algonotes}

%% ========== 空间OLS回归 ==========
\basicalgo{空间OLS回归}{spatial_ols}

\begin{algointro}
先做普通最小二乘回归（OLS），然后\textbf{检验残差有没有空间自相关}。如果有，说明普通回归不可靠，需要改用空间模型。它是空间分析的"基准诊断"步骤。
\end{algointro}

\begin{algoscene}
\begin{itemize}
    \item 任何空间回归分析的\textbf{第一步}：先做OLS基准，再诊断
    \item 判断"普通回归的结论可不可信"
    \item 为选择SAR还是SEM提供依据（LM检验）
\end{itemize}
\end{algoscene}

\begin{algoprinciple}
普通回归假设 $y = X\beta + \varepsilon$，且 $\varepsilon$ 之间互相独立。但现实中相邻地区的残差往往相关（比如一个政策同时影响相邻几个区县）。

步骤：
\begin{enumerate}
    \item 先跑普通OLS，得到系数和残差
    \item 对残差做\textbf{Moran's I检验}：如果残差Moran's I显著，说明残差有空间聚集→OLS假设被破坏
    \item 做\textbf{LM检验}：LM-error显著→选SEM；LM-lag显著→选SAR；都显著→看稳健形式
\end{enumerate}

就像体检：先做基础检查（OLS），再做进阶筛查（空间诊断），决定要不要进一步治疗（空间模型）。
\end{algoprinciple}

\begin{algosposs}
在\textbf{Dabbit AI SPSS工具箱}中操作：

\begin{enumerate}
    \item 上传区域数据（因变量、自变量、坐标）
    \item 指定因变量（如地均GDP）和自变量（投资强度、研发强度等）
    \item 权重方法选KNN（k=4），距离度量haversine，行标准化
    \item 稳健标准误选"否"
    \item 运行后先看OLS系数表，再看空间依赖诊断表
\end{enumerate}
\end{algosposs}

\begin{algoresult}
\begin{itemize}
    \item \textbf{OLS系数}：如投资强度系数2.09（$p<0.001$），研发强度1.51（$p<0.001$）
    \item \textbf{$R^2=0.63$}：模型解释力
    \item \textbf{残差Moran's I=0.31（$p<0.001$）}：残差显著空间自相关→OLS不可靠
    \item \textbf{LM检验}：LM-Error显著、稳健LM-Error显著→建议用空间误差模型SEM
\end{itemize}
\end{algoresult}

\begin{algonotes}
\textbf{什么时候用？}空间分析的起点，必须先做。

\textbf{什么时候不用？}如果残差Moran's I不显著，说明没有空间效应，OLS即可，不需要后续空间模型。

\textbf{常见坑}：跳过空间诊断直接做空间模型——应先证明"需要"空间模型。
\end{algonotes}

%% ========== Moran指数（空间自相关） ==========
\basicalgo{{\starmark}Moran指数（空间自相关）}{moran_i}

\begin{algointro}
Moran's I是衡量\textbf{空间聚集程度}的核心指标。它回答一个问题：高值区域是不是和高值区域相邻？低值和低值相邻？还是完全随机散布？比如PM2.5是不是连片污染？房价是不是核心区高、外围低？
\end{algointro}

\begin{algoscene}
\begin{itemize}
    \item \textbf{环境科学}：PM2.5污染是否连片聚集？
    \item \textbf{城市经济}：高GDP区县是否空间聚集？
    \item \textbf{房地产}：高价小区是否扎堆？
    \item 任何"地理上相近的观测值是否相似"的问题
\end{itemize}
\end{algoscene}

\begin{algoprinciple}
Moran's I的公式本质是\textbf{变量值与其空间滞后的相关系数}：

\[ I = \frac{n}{\sum_i\sum_j w_{ij}} \cdot \frac{\sum_i\sum_j w_{ij}(x_i-\bar{x})(x_j-\bar{x})}{\sum_i(x_i-\bar{x})^2} \]

\begin{itemize}
    \item $I>0$且显著：\textbf{正空间自相关}——高值邻高值、低值邻低值（连片聚集）
    \item $I<0$且显著：\textbf{负空间自相关}——高值邻低值（高低交错）
    \item $I\approx 0$：空间随机分布，没有聚集
\end{itemize}

\textbf{Moran散点图}：横轴是变量值（标准化），纵轴是邻居均值。一、三象限=高-高/低-低聚集；二、四象限=高-低/低-高离群。

就像看地图：如果红色（高值）都连成片、蓝色（低值）都连成片，Moran's I就大；如果红蓝交错，I就接近0。
\end{algoprinciple}

\begin{algosposs}
SPSS原生不支持Moran指数，在\textbf{Dabbit AI SPSS工具箱}中操作：

\begin{enumerate}
    \item 上传含坐标的数据
    \item 变量列选要检验的变量（如PM2\_5年均浓度）
    \item 权重构建方式选"k近邻"，k=4
    \item 行标准化选"是"
    \item 置换检验次数设为999，显著性水平0.05
    \item 固定随机种子42（可复现）
    \item 点击运行
\end{enumerate}
\end{algosposs}

\begin{algoresult}
\begin{itemize}
    \item \textbf{Moran's I=0.84}（期望$E[I]=-0.007$）：远大于期望
    \item \textbf{z值=15.2，置换p=0.002 $<0.05$}：显著正空间自相关
    \item \textbf{象限占比}：高-高38\%、低-低54\%——大部分是同类聚集
    \item \textbf{结论}：PM2.5浓度呈连片聚集分布，不是随机散布
\end{itemize}

\textbf{注意}：全局Moran's I只告诉你"整体有没有聚集"，不能定位\textbf{哪里}是热点（那需要局部Moran's I/LISA）。
\end{algoresult}

\begin{algonotes}
\textbf{什么时候用？}数据有地理位置，想判断"是不是空间聚集"时。

\textbf{什么时候不用？}
\begin{itemize}
    \item 数据没有空间坐标——无法算Moran's I
    \item 已经在回归残差上做过Moran检验——那是诊断残差，这里是检验因变量本身
\end{itemize}

\textbf{常见坑}
\begin{itemize}
    \item \textbf{坑1}：把Moran's I当相关系数解读——I的取值范围不是[-1,1]，而是近似$[-1/(n-1), 1]$
    \item \textbf{坑2}：全局I显著就说"所有区域都聚集"——全局I只说整体趋势，个别区域可能是离群点
\end{itemize}

\textbf{和空间回归的关系}：Moran's I是"要不要做空间模型"的诊断；空间回归是"做了之后怎么解释效应"。
\end{algonotes}

%% ========== 空间滞后模型（SLM/SAR） ==========
\basicalgo{{\starmark}空间滞后模型（SLM/SAR）}{sar_model}

\begin{algointro}
空间滞后模型（SAR/SLM）在普通回归中加入\textbf{邻居因变量} $Wy$——不仅看本地自变量，还看邻居的因变量对本地的影响。比如：一个小区涨价会带动周边小区跟涨。
\end{algointro}

\begin{algoscene}
\begin{itemize}
    \item \textbf{房地产估价}：小区挂牌价受周边小区价格影响
    \item \textbf{区域经济}：一个城市的GDP受相邻城市增长带动
    \item \textbf{创新扩散}：一个地区的创新产出受邻近地区知识溢出影响
\end{itemize}
\end{algoscene}

\begin{algoprinciple}
普通回归：$y = X\beta + \varepsilon$

空间滞后模型：$y = \rho W y + X\beta + \varepsilon$

\begin{itemize}
    \item $Wy$是"邻居因变量的加权平均"
    \item $\rho$（空间滞后系数）：邻居每涨1单位，本地涨$\rho$单位
    \item $\rho>0$显著：存在\textbf{正向空间溢出}（邻居好我也好）
\end{itemize}

关键：因为$Wy$和$y$互相影响（你影响我、我影响你），形成了\textbf{空间反馈循环}，所以不能直接用OLS估计，要用极大似然（ML）。

效应分解：
\begin{itemize}
    \item \textbf{直接效应}：本地自变量对本地因变量的影响
    \item \textbf{间接效应（溢出）}：本地自变量变化通过空间网络传导到邻居
    \item \textbf{总效应}=直接+间接
\end{itemize}
\end{algoprinciple}

\begin{algosposs}
在\textbf{Dabbit AI SPSS工具箱}中操作：

\begin{enumerate}
    \item 上传含坐标和变量的数据
    \item 因变量选"avg\_price"，自变量选"metro\_km, school\_score, building\_age"
    \item 坐标选lon/lat，权重类型knn近邻，k=4，行标准化
    \item 估计方法选极大似然ML
    \item 点击运行，输出系数表和效应分解表
\end{enumerate}
\end{algosposs}

\begin{algoresult}
\begin{itemize}
    \item \textbf{$\rho=0.563$（$p<0.001$）}：显著正向空间溢出——邻居均价每涨1000元，本地涨约563元
    \item \textbf{LR检验}：$\rho=0$被拒绝→SAR比OLS显著更优
    \item \textbf{效应分解}：如距地铁距离的直接效应-0.53、间接效应-0.58、总效应-1.11——间接占比52\%！
    \item \textbf{间接占比}：超过50\%说明空间溢出效应非常强，不能忽略
\end{itemize}
\end{algoresult}

\begin{algonotes}
\textbf{什么时候用？}Moran's I显著，且LM-lag检验显著时。

\textbf{什么时候不用？}
\begin{itemize}
    \item LM-error显著而LM-lag不显著→用SEM
    \item Moran's I不显著→用普通OLS即可
\end{itemize}

\textbf{常见坑}
\begin{itemize}
    \item \textbf{坑1}：把SAR系数直接当OLS系数解读——SAR中自变量系数只是直接效应，总效应要算间接效应
    \item \textbf{坑2}：忽略$Wy$的内生性——$Wy$和$y$互为因果，OLS估计有偏，必须用ML或IV
\end{itemize}

\textbf{和SEM的区别}：SAR是\textbf{因变量}之间的空间溢出（邻居y影响本地y）；SEM是\textbf{误差项}之间的空间相关（未观测因素在相邻地区间传导）。
\end{algonotes}

\section{进阶级算法速览}

%% ========== 空间误差模型（SEM） ==========
\advancedalgo{空间误差模型（SEM）}{spatial_error_model}

\begin{algobrief}
\textbf{一句话}：空间依赖体现在\textbf{误差项}中——未观测的扰动在相邻地区间传导，而非因变量本身有溢出。

\textbf{美赛场景区}：房价分析中，相邻片区共享未观测的地段口碑、学区预期等因素，导致残差空间相关。

\textbf{核心原理}：$y=X\beta+u$，$u=\lambda Wu+\varepsilon$。$\lambda$是空间误差系数。LM-error显著而LM-lag不显著时选SEM。

\textbf{操作要点}：KNN权重k=4，ML估计，行标准化。

\textbf{结果关键}：$\lambda=0.537$（$p<0.001$）显著，说明误差存在空间扩散；系数与OLS接近但标准误更准。

\textbf{注意}：SEM中自变量没有空间溢出项，其效应就是直接效应，不需要分解。
\end{algobrief}

%% ========== 空间杜宾模型（SDM） ==========
\advancedalgo{空间杜宾模型（SDM）}{sdm_model}

\begin{algobrief}
\textbf{一句话}：最全面的空间模型——同时包含因变量空间滞后$Wy$和自变量空间滞后$WX$，既看邻居y也看邻居x。

\textbf{美赛场景区}：区域创新研究——本地研发投入影响本地创新，邻居研发投入也通过知识溢出影响本地。

\textbf{核心原理}：$y=\rho Wy+X\beta+WX\theta+\varepsilon$。可以检验能否退化为SAR（$\theta=0$）或SEM。SDM是最一般化的设定。

\textbf{操作要点}：ML估计，wx列选"全部解释变量"，行标准化。

\textbf{结果关键}：$\rho=0.45$显著；研发投入直接效应1.12、间接效应0.81——溢出效应很大。LR检验SDM显著优于SAR和SEM。

\textbf{注意}：SDM最全面但参数最多，小样本时估计不稳定；大样本首选SDM。
\end{algobrief}

%% ========== 空间滞后误差模型（SAC） ==========
\advancedalgo{空间滞后误差模型（SAC）}{sac_model}

\begin{algobrief}
\textbf{一句话}：同时包含空间滞后$\rho Wy$和空间误差$\lambda Wu$，SAR和SEM的合体。

\textbf{美赛场景区}：住宅价格中既有邻居房价的实质溢出（$\rho$），又有未观测区位因素的误差相关（$\lambda$）。

\textbf{核心原理}：$y=\rho Wy+X\beta+u$，$u=\lambda Wu+\varepsilon$。两个空间参数都显著说明两种机制都存在。

\textbf{操作要点}：ML估计，KNN权重，行标准化。

\textbf{结果关键}：$\rho=0.479$和$\lambda=0.401$都显著→两种空间机制同时存在。

\textbf{注意}：SAC参数多，对样本量要求高；小样本时$\rho$和$\lambda$可能难以区分。
\end{algobrief}

%% ========== 空间面板模型 ==========
\advancedalgo{空间面板模型}{spatial_panel}

\begin{algobrief}
\textbf{一句话}：空间模型+面板数据（多地区×多年份），同时控制时间不变的地区固定效应和空间溢出。

\textbf{美赛场景区}：15个城市10年的经济增长面板，研究数字经济对GDP的拉动及跨城溢出。

\textbf{核心原理}：在SAR/SEM/SDM基础上加入个体固定效应（吸收各城市不随时间变的禀赋差异），用组内去均值后做ML估计。

\textbf{操作要点}：effects选fixed，model type选sar，KNN权重k=4。

\textbf{结果关键}：$\rho=0.424$显著（城市间经济联动）；数字人才直接效应15.5、间接溢出10.6——近40\%效应通过邻国外溢。

\textbf{注意}：面板数据要有足够时间维度（$T>5$），固定效应吸收了不随时间变的因素。
\end{algobrief}

%% ========== 空间杜宾误差模型（SDEM） ==========
\advancedalgo{空间杜宾误差模型（SDEM）}{sdem_model}

\begin{algobrief}
\textbf{一句话}：包含自变量空间滞后$WX$和误差空间相关$\lambda$，但\textbf{不包含}因变量空间滞后$Wy$。

\textbf{美赛场景区}：创新产出研究——邻居研发投入有溢出（WX），但误差项也有空间关联（$\lambda$）。

\textbf{核心原理}：$y=X\beta+WX\theta+u$，$u=\lambda Wu+\varepsilon$。与SDM的区别：SDEM没有$Wy$，所以效应分解简单——$\beta$是直接效应，$\theta$就是间接效应，不需要矩阵求逆。

\textbf{操作要点}：KNN权重k=4，ML估计，行标准化。

\textbf{结果关键}：$\lambda=0.502$显著（误差空间相关）；研发投入直接效应1.10、间接效应0.78。

\textbf{注意}：SDEM比SDM简单（不需要$(I-\rho W)^{-1}$），如果LM-lag不显著但LM-ERR显著，SDEM是好选择。
\end{algobrief}

%% ========== 自变量空间滞后模型（SLX） ==========
\advancedalgo{自变量空间滞后模型（SLX）}{slx_model}

\begin{algobrief}
\textbf{一句话}：最简单的空间模型——只加邻居自变量$WX$，不加$Wy$也不加误差项，用普通OLS就能估计。

\textbf{美赛场景区}：便利店研究——邻居门店的营销投入是否影响本店营业额（客流共享/分流）。

\textbf{核心原理}：$y=X\beta+WX\theta+\varepsilon$。$\beta$=本地自变量效应（直接），$\theta$=邻居自变量效应（间接溢出）。因为没有$Wy$内生性，直接OLS即可。

\textbf{操作要点}：knn k=4，lag mode=all，稳健标准误，行标准化。

\textbf{结果关键}：营销投入直接效应1.56、间接效应1.08——邻居营销也能带动本店（品牌集聚效应）。

\textbf{注意}：SLX是入门级空间模型，如果残差Moran's I仍显著，说明还需要加误差项（SDEM）。
\end{algobrief}
